\subsection{Complete residue-selection policy}
\label{sec:identity_details}
Only the first PDB model is analyzed. Standard amino acids and modified residues recognized by Biopython's extended protein-residue mapping contribute one selected C-alpha atom. Selenomethionine is mapped to methionine in the observed sequence; SEC, PYL, and UNK are represented as X. Twelve additional components classified as D-peptide linking by the Chemical Component Dictionary \citep{ccd} are retained geometrically and represented as X, without asserting an L-amino-acid parent identity: DIV, DAL, DPN, DVA, DCY, DPR, DGL, DTR, DAR, DLE, DAS, and DLY. The three-letter identities remain available in the residue audit. The installed mapping and component-specific source links are recorded in the configuration.

Nonpolymer components CA, PLM, and MPT are explicitly excluded, even where the supplied coordinate file labels their records ATOM. An atom name of CA alone does not establish a protein C-alpha. The final audit retains 78,400 residues in 364 structures and excludes 19 nonpolymer records from the historical 78,419-row collection: 16 calcium ions, one PLM record, and two MPT records. Relative to an initial standard-residue audit, this policy restores 177 modified or D-peptide residues before fitting any model. Geometry includes these retained coordinates; unrecognized components receive no generic carbon-atom fallback.

Among alternate C-alpha locations, choose the highest occupancy, breaking ties by blank alternate-location code, then A, then lexical order. Duplicate records with the same residue and alternate-location key are treated as an error. The element field must be carbon or blank; a conflicting explicit element is rejected.

Nonfinite coordinates, occupancies, or B-factors and negative occupancies or B-factors are rejected. Structures with fewer than three retained C-alpha atoms or constant retained B-factors are excluded from correlation-based prediction. Each selected row preserves the structure, model, chain, residue number, insertion code, atom selection, and source hash. Coincident coordinates require an explicit geometry check before restrictions dividing by distance can be computed. A preprocessing exclusion and a numerical feature-generation failure are separate audit categories.

Normalized B-factor is computed over the complete retained protein entry, with population standard deviation, and is never a feature. Legacy annotation tables are excluded from the primary feature matrix because a supplied residue index alone does not establish a verified chain/residue mapping. Source annotation correlations from historical analyses therefore are not claimed as validation of the new model. The present primary target and all predictors come from the audited structure rows.
